\section{Implementation and Experiments}\label{sec:experiments}

\subsection{Reference Implementation}

The correspondence is implemented in the Operon framework (\texttt{operon-ai}, MIT license). Key components:

\begin{itemize}[leftmargin=*]
\item \textbf{Architecture extraction}: \texttt{extract\_architecture(organism)} produces the $(G, \mathrm{Know}, \Phi)$ triple from any \texttt{SkillOrganism}.
\item \textbf{Compiler functors}: Five \texttt{CompilerFunctor} instances with automatic preservation verification.
\item \textbf{Certificate framework}: Self-verifying certificates with \texttt{certify()} and \texttt{verify()} methods on certifiable components (ATP budget, quorum sensing, mTOR scaling).
\item \textbf{Atomic skills catalog}: \texttt{seed\_library\_from\_atomic\_skills()} registers the five skills with correct topology mapping.
\item \textbf{LangGraph compiler}: \texttt{organism\_to\_langgraph()} creates one LangGraph node per stage, each calling \texttt{run\_single\_stage()}.
\end{itemize}

\subsection{Escalation Experiment}

To test whether structural guarantees are genuinely harness-level, we designed an experiment that changes the model while keeping the harness constant.

\paragraph{Setup.} A single-stage code review organism with:
\begin{itemize}[leftmargin=*]
\item \textbf{Fast model}: Phi-3 Mini (3.8B parameters) via Ollama
\item \textbf{Deep model}: Gemma 4 (27B MoE, 4B active) via Ollama
\item \textbf{VerifierComponent}: Rubric-based quality evaluation using Gemma~4 as judge
\item \textbf{WatcherComponent}: Convergence detection with quality-based escalation
\item \textbf{Quality threshold}: 0.6 (below this, the watcher fires ESCALATE)
\end{itemize}

The task is \texttt{hard\_par\_08}: a code review containing four subtle bugs (off-by-one pagination, TOCTOU cache race, float precision in financial math, \texttt{except Exception} swallowing \texttt{KeyboardInterrupt}).

\paragraph{Results.}

\begin{center}
\small
\begin{tabular}{@{}lccl@{}}
\toprule
\textbf{Phase} & \textbf{Model} & \textbf{Quality} & \textbf{Action} \\
\midrule
Initial execution & Phi-3 Mini (fast) & 0.50 & EXECUTE \\
Verifier evaluation & Gemma 4 (judge) & --- & Score = 0.50 \\
Watcher decision & --- & $0.50 < 0.60$ & ESCALATE \\
Escalated execution & Gemma 4 (deep) & --- & EXECUTE \\
\bottomrule
\end{tabular}
\end{center}

The escalation fires correctly: Phi-3 Mini produces a review that the judge scores below threshold, the watcher escalates to Gemma~4, and the deep model re-executes the stage. The structural guarantee (quality-based escalation) operates independently of which model executes---it is a property of the harness ($\mathrm{Know}$), not the model ($\Phi$).

\paragraph{Discrimination.} Independent validation of the task confirms discrimination between weak and strong models: Phi-3 Mini scores 0.72 (identifies bug classes but not concrete failure scenarios), Gemma~4 scores 1.00 (explains each failure scenario in detail). Quality delta = 0.28, exceeding the 0.20 threshold.

\subsection{Certificate Preservation Measurement}

We compile organisms with multiple certificates (ATP priority gating + quorum sensing no-false-activation + mTOR no-oscillation) across all four non-LangGraph compilers and measure preservation:

\begin{center}
\small
\begin{tabular}{@{}lcc@{}}
\toprule
\textbf{Compiler} & \textbf{Certificates Preserved} & \textbf{Verified} \\
\midrule
Swarms & 3/3 (100\%) & 3/3 \\
DeerFlow & 3/3 (100\%) & 3/3 \\
Ralph & 3/3 (100\%) & 3/3 \\
Scion & 3/3 (100\%) & 3/3 \\
\bottomrule
\end{tabular}
\end{center}

Full certificate identity (theorem + parameters + source) is preserved across all four measured compilers, confirming Proposition~5.1 empirically. The LangGraph compiler preserves certificates via per-stage execution of \texttt{run\_single\_stage()} (Section~\ref{sec:preservation}).

\subsection{SWE-bench-lite: An Inconclusive Run Dominated by Patch-Apply Failures}
\label{sec:swebench-phase2}

To probe whether the categorical harness delivers end-to-end value beyond certificate preservation, we attempted an apples-to-apples comparison on SWE-bench-lite: 10 real Python bug-fix instances from \texttt{astropy} and \texttt{django}, with patches applied inside the official SWE-bench Docker harness and evaluated via \texttt{FAIL\_TO\_PASS} + \texttt{PASS\_TO\_PASS} test suites. The run is reported here for transparency, but --- as we show below --- it is confounded and cannot support a capability conclusion.

\paragraph{Setup.} Gemma~4 via Ollama (tag \texttt{gemma4:latest}, digest \texttt{c6eb396dbd59}, blob \texttt{sha256:4c27e0f5}, architecture \texttt{gemma4}, 8.0B parameters, Q4\_K\_M quantization, 131{,}072 context). The immutable digest is recorded in \texttt{eval/results/swebench\_phase2.json} so the exact weights can be identified later. Three conditions:
\begin{itemize}[leftmargin=*]
\item \textbf{baseline}: direct single prompt, request unified diff.
\item \textbf{organism}: three-stage \texttt{SkillOrganism} (\textit{localize} $\to$ \textit{edit} $\to$ \textit{verify}).
\item \textbf{langgraph}: same organism, compiled via \texttt{organism\_to\_langgraph()} and executed as a LangGraph \texttt{StateGraph}.
\end{itemize}

\paragraph{Results (harness outcome distribution, 10 instances per condition).}

\begin{center}
\small
\begin{tabular}{@{}lcccccc@{}}
\toprule
\textbf{Condition} & \textbf{Resolved} & \textbf{Unresolved} & \textbf{Error} & \textbf{Empty patch} & \textbf{Evaluated} & \textbf{Mean latency} \\
\midrule
baseline    & 0 & 1 & 9 & 0 & 1/10 & 44~s  \\
organism    & 0 & 1 & 5 & 4 & 1/10 & 88~s  \\
langgraph   & 0 & 0 & 6 & 4 & 0/10 & 90~s  \\
\bottomrule
\end{tabular}
\end{center}

Categories follow the SWE-bench harness: \texttt{resolved} (patch applied and all target tests pass), \texttt{unresolved} (patch applied, tests fail), \texttt{error} (patch rejected by \texttt{git apply}, timeout, or infrastructure failure), \texttt{empty\_patch} (no diff-shaped output from the model). \textbf{Evaluated} counts only instances that reached a pass/fail harness verdict, i.e.\ \texttt{resolved} $+$ \texttt{unresolved}.

\paragraph{Interpretation (confounded; no capability conclusion).} The run does not isolate model capability from infrastructure. Across all three conditions, only 2 of 30 submissions (baseline: 1, organism: 1, langgraph: 0) reached an actual test-suite verdict; the remaining 28 were rejected at the \texttt{git apply} stage (\texttt{error}, 20) or never produced a diff-shaped output at all (\texttt{empty\_patch}, 8). With \textbf{evaluated} in $\{0, 1\}$ per condition, the sample size is too small to distinguish model weakness from upstream patch-format/pathing issues. We therefore do \emph{not} claim this run demonstrates a capability ceiling.

Two observations are still defensible from the data we have:

\begin{enumerate}[leftmargin=*]
\item \textbf{The organism and LangGraph pipelines more often failed to produce a diff-shaped output} (4 \texttt{empty\_patch} each vs.\ 0 for baseline). This is a measurement about extractability of the three-stage pipeline's output, not about patch correctness: baseline produced an extractable diff for every instance yet 9/10 were rejected by \texttt{git apply} anyway, so extractability is not a proxy for patch quality.
\item \textbf{Latency roughly doubled under the pipelines} (88--90~s vs.\ 44~s mean) with no compensating gain in evaluated instances. Cost is paid regardless of whether the output is usable.
\end{enumerate}

The Phase~1 LLM-judge scores reported earlier (baseline~$=0.808$, organism~$=0.776$) are plausibility judgments over the produced text; they do not measure whether a patch applies to the target repository, which is where most of this run failed. That is an argument for using harness-grounded pass/fail rather than judge plausibility for end-task claims, but it does not by itself rank the three conditions.

\paragraph{What this run does and does not say.} It does not say that categorical compilation or multi-stage decomposition fail on SWE-bench-lite; with 0--1 evaluated instances per condition, the experiment is not powered to answer that. What it does say is that (i) the current prompt-to-unified-diff pipeline produces patches that are almost never applicable to the target repos, and (ii) before decomposition vs.\ direct prompting can be compared on this benchmark, the upstream patch-application failure rate has to be brought down. The harness-level \emph{structural} guarantees reported in Sections~\ref{sec:preservation} and in the escalation experiment are unaffected by this negative run: they are $\mathrm{Know}$-level properties verified independently of SWE-bench outcomes.

This motivates Section~\ref{sec:conclusion}: the demonstrated transfer value of categorical harness engineering in this paper is in \emph{preservation} and \emph{primitive reuse}, not in task-resolution gains on SWE-bench-lite --- the latter remains open pending a run where patches reliably apply.
